% Copyright 2026 Torbjörn E. M. Nordling
% SPDX-License-Identifier: Apache-2.0
% frame-run-pi-e1.tex — Run the experiment on the Pi (1/2): connect, E1 (21:15–21:45)
% Hand-out level: shown briefly, Sam demonstrates the same commands next.
% Commands for robot R5, left wheel, motor m9: Sam's instructions for running
% the tests on the Pi (approved by the user, 2026-11-12). Added after a
% usability test with a novice student. The trailing backslash continues a
% shell command on the next line, so each block can be typed as shown.
\begin{frame}[fragile]{Run the experiment on the Pi (1/2)}{Example: robot R5, left wheel, motor m9 --- type your own}
  \scriptsize
  \textbf{Connect} from a terminal (Git Bash on Windows); the TAs give you the user, address, and password in class:
\begin{verbatim}
ssh <user>@<address>
sudo systemctl stop cherry   # stop the robot's driver; it holds the motors
cd ~/experiments
tmux new -A -s run           # the run continues if the connection drops
\end{verbatim}
  \vspace{-4pt}
  \textbf{E1, one video, both directions.} Start the video first, LED in view; then:
\begin{verbatim}
python3 motor_trial.py A --robot R5 --side L --motor m9 \
  --camera-fps 240 --direction forward \
  --start 0.10 --top 0.36 --step 0.02 --hold 2 --repeat 1
\end{verbatim}
  \vspace{-4pt}
  When it prints \texttt{Done (completed)}, run the same command with \verb|--direction backward|; then stop the video. Name it:\\
  \texttt{R5-L-m9\_A\_sweep\_dense\_forward-backward\_010to036by020\_rep01\_<YYYYMMDD>}
\end{frame}
